\documentclass[11pt,a4paper]{amsart}
\usepackage{amsmath,amssymb,amsthm,booktabs}
\usepackage[margin=1in]{geometry}
\usepackage{hyperref}
\newtheorem{proposition}{Proposition}[section]
\newtheorem{remark}[proposition]{Remark}
\newcommand{\Q}{\mathbb Q}
\newcommand{\C}{\mathbb C}
\newcommand{\R}{\mathbb R}
\newcommand{\Z}{\mathbb Z}
\newcommand{\SU}{\mathrm{SU}}
\newcommand{\SL}{\mathrm{SL}}
\newcommand{\Spin}{\mathrm{Spin}}
\newcommand{\Sym}{\mathrm{Sym}}
\title{Exact Representation-Theoretic Checks for Binary Icosahedral Constructions}
\author{Tyler}
\address{Independent researcher}
\date{7 October 2026 --- corrected revision}
\begin{document}
\begin{abstract}
We report exact finite calculations for the binary icosahedral group
$2I=\SL(2,5)$, selected maps into exceptional Lie groups, and
$\SL(2,\Z/70)$. We separate verified representation-theoretic facts from
claims about compactifications and physical chirality. The character and
Frobenius--Schur table, an eight-dimensional spinor-branching search and
low-degree tensor invariants admit exact certificates. Several stronger
exclusions in the original formulation do not follow: outer tensor products
can carry mixed character fields, self-duality alone does not force a zero
Dirac index, and selected examples do not classify all compactifications.
The supplied scripts verify the explicitly stated finite scope, not a
universal Standard-Model impossibility theorem.
\end{abstract}
\maketitle

\section{Scope and verification conventions}
The computations use exact integer pairs $a+b\varphi$, where
$\varphi^2=\varphi+1$, for the binary-icosahedral group. Quaternion
coordinates are $(a+b\varphi)/2$. Finite class algebras and cyclotomic
polynomials are handled by exact rational and algebraic arithmetic.
A script failure is distinct from a disproved conjecture; uncomputed
catalog or physics claims are recorded as unverified. The programs are
executable certificates, not a proof-assistant formalization. This revision
supersedes the mathematical claims of the historical preprint; the original
source and PDF remain archived under their original names. Verification
certificates apply only to the corrected scope stated here.

The internal space $S^3/2I\times T^2$ has real dimension five. Adding
four-dimensional spacetime gives dimension nine. A ten-dimensional
heterotic interpretation therefore requires additional explicitly specified
structure; this product alone does not supply a six-dimensional internal
space. No universal classification of such additional structure is attempted.

\section{Characters and arithmetic}
\begin{proposition}
For a finite group $G$, all ordinary character values are algebraic
integers in $\Q(\zeta_{\exp G})$. For $G=2I$, the exponent is $60$,
and $\sin(\pi/14)\notin\Q(\zeta_{60})$.
Moreover, $\sin(\pi/14)$ is not an ordinary character value of any finite group.
\end{proposition}
\begin{proof}
The eigenvalues of $g$ in a finite-dimensional representation are roots of
unity whose orders divide $\exp G$, so their sum has the asserted properties.
Set $z=e^{\pi i/7}$. Since $z^7=-1$,
$s=(z^3-z^4)/2=\cos(3\pi/7)=\sin(\pi/14)$.
Exact reduction modulo $\Phi_{14}(z)$ gives
$8s^3-4s^2-4s+1=0$. The polynomial is irreducible over $\Q$;
its possible rational roots $\pm1,\pm1/2,\pm1/4,\pm1/8$ are not roots.
Thus $[\Q(s):\Q]=3$, which cannot divide
$[\Q(\zeta_{60}):\Q]=\phi(60)=16$.
The monic minimal polynomial
$x^3-\tfrac12x^2-\tfrac12x+\tfrac18$ is not in $\Z[x]$, so $s$ is not an
algebraic integer. This proves the literal-character assertion, not a
statement about normalized observables or mixing angles.
\end{proof}

The exact character field of $2I$ is $\Q(\sqrt5)$; the exponent field above
is an upper bound. In particular, arithmetic must be applied to the
specific data and permitted operations, not simply to a level label.

\begin{proposition}
The nine irreducible characters of $2I$ have the following dimensions and
Frobenius--Schur indicators. Precisely the quaternionic irreducibles are
faithful.
\end{proposition}
\begin{center}
\begin{tabular}{lrrrrrrrrr}
\toprule
Label & $1$ & $2_a$ & $2_b$ & $3_a$ & $3_b$ & $4_R$ & $4_H$ & $5$ & $6$\\
Dimension & 1&2&2&3&3&4&4&5&6\\
Indicator & 1&-1&-1&1&1&1&-1&1&-1\\
\bottomrule
\end{tabular}
\end{center}
\begin{proof}
The 120 exact unit quaternions give nine conjugacy classes, with sizes
\[1,1,12,12,12,12,20,20,30.\] Starting with the defining doublet,
$\chi_{\Sym^{n+1}}=\chi_2\chi_{\Sym^n}-\chi_{\Sym^{n-1}}$ supplies the
symmetric powers. The second doublet and triplet are their Galois
conjugates; $4_R=2_a\otimes2_b$. These nine actual characters have exact
inner products $\delta_{ij}$ and dimension-square sum $120$.
Evaluating $|G|^{-1}\sum_g\chi(g^2)$ gives the displayed indicators.
Kernels are checked directly using $\chi(g)=\chi(1)$, and
$\sum_\chi\nu(\chi)\chi(1)=2=\#\{g:g^2=1\}$.
\end{proof}

\begin{remark}[What self-duality does not imply]
This table does not prove that every equivariant Dirac index vanishes.
An equivariant index is $[\ker D^+]-[\ker D^-]\in R(G)$, and a nonzero
virtual representation can be self-dual. An antilinear map commuting with
the chirality grading preserves, rather than exchanges, the two graded
spaces. Vanishing requires additional structure pairing opposite-chirality
kernels, compatible with the complete geometry and bundle twists. These
conditions are not established by a character table alone.
\end{remark}

\section{Eight-dimensional spinors and invariants}
The real irreducible pieces have dimensions $d$ for real type and $2d$
for quaternionic type. There are 25 multisets of total real dimension eight,
12 faithful. The programs construct lifts using
$\Spin(4)=\SU(2)\times\SU(2)$ and the lift of the realified
$\Sym^3\C^2$ representation, not independent choices of angles on classes.
For two four-dimensional blocks with half-spinors $(A,B)$ and $(C,D)$,
\[
 V=AB+CD,\qquad S^+=AC+BD,\qquad S^-=AD+BC.
\]
The realified $4_H$ has unordered half-spinor pair
$\{3\cdot1+5,\,2\cdot4_H\}$. Exact decomposition shows that it is the
unique faithful real-eight module with three singlets in one half-spinor.
This counts invariant spinors of the specified representations, not
physical generations in an unspecified compactification.

For the defining doublet and triplet,
$2_a\otimes2_a=1+3_a$, so this tensor product does contain a singlet.
Using the harmonic decomposition
$\Sym^n(3_a)=V_n+V_{n-2}+\cdots$ gives
\[
\dim(2_a\otimes2_a\otimes\Sym^n3_a)^{2I}
=1,1,1,1,1,2,3,3,3,4,5,6,7
\quad(0\leq n\leq12).
\]
The doublet Molien coefficients agree through degree 60 with
$(1+t^{30})/((1-t^{12})(1-t^{20}))$.

\section{Selected centralizers}
Write $V_j$ for the spin-$j$ representation of $\SU(2)$.
For a principal $\SU(2)$, the adjoint is the sum of $V_j$ with $j$ the
Lie-algebra exponents, without added Cartan singlets. The programs recover
these branchings from exact root and weight orbits and cross-check the
$E_8$ branching against standard tables \cite{SageBranching}.
\begin{center}
\begin{tabular}{lr}
\toprule
Specified restricted adjoint & Dimension of $2I$ invariants\\
\midrule
Principal $\SU(2)$ in $E_8$ & 0\\
Principal $\SU(2)$ in $E_6$ & 0\\
Principal $\SU(2)$ in $E_6\subset E_8$ & 14\\
Regular $\SU(2)$ in $\SU(6)\times\SU(2)\subset E_6\subset E_8$ & 133\\
Regular $\SU(2)$ in $E_7\times\SU(2)\subset E_8$ & 133\\
\bottomrule
\end{tabular}
\end{center}
For the $E_6\times A_2$ chain, one must retain the mixed terms:
\[
248=(78,1)+(1,8)+(27,3)+(\overline{27},\overline3).
\]
The principal restriction of $27$ is $V_0+V_4+V_8$, with one $2I$ invariant;
thus the full invariant dimension is $0+8+6=14$.
For the regular $\SU(2)$, the corresponding invariant dimensions in $78$
and $27$ are 35 and 15, giving $35+8+6\cdot15=133$.
The principal maps in $E_6$ and $E_8$ kill the central involution of $2I$;
their finite image is $A_5$, not an injectively embedded $2I$.

For the particular block embedding
$\R^{16}=(\Sym^3\C^2)_{\R}\oplus\R^8$, the fixed dimensions in the
$D_8$ adjoint and half-spinor are respectively 31 and 24. There is an
additional identification argument, rather than a dimension-only guess:
simultaneous triality on the two $D_4$ blocks preserves the 240 roots of
$E_8$ and takes the torus weights
$(3,3,1,1,0,0,0,0)$ to $(4,2,0,0,0,0,0,0)$.
The vector then splits as $V_2\oplus\R^{11}$, exhibiting a commuting
$\mathfrak{so}(11)$ subalgebra. Under $\Spin(5)\times\Spin(11)$,
\[
248=(10,1)+(1,55)+(5,11)+(4,32),
\]
whose first-factor restrictions are $V_1+V_3$, $V_0$, $V_2$ and $V_{3/2}$.
The nontrivial summands have no $2I$ invariants. Thus the exhibited
$\mathfrak{so}(11)$ exhausts the 55 fixed dimensions: the fixed Lie algebra
is $B_5$, of rank five, not $F_4\oplus\mathfrak u(1)^3$.
The global centralizer group and its components are not identified here.

The standard branching
$248=(52,1)+(1,14)+(26,7)$ under $F_4\times G_2$ contains no trivial
summand, hence its Lie-algebra centralizer is zero \cite{SageBranching}.
Consequently no nontrivial \emph{connected} holonomy subgroup can commute
with this entire pair. A claim about the full group centralizer, disconnected
holonomy, or an asymmetric orbifold requires additional arguments.
$F_4\times G_2$ is a special, non-regular subalgebra here; it is not a
symmetric pair or a Dynkin-diagram folding of $E_8$.

\section{Congruence groups and counterexamples to broad exclusions}
CRT gives
$\SL(2,\Z/70)\cong S_3\times2I\times\SL(2,7)$.
The group has order $241920$, exponent $840$, and 297 irreducibles; their
Frobenius--Schur counts are 96 real, 93 quaternionic and 108 complex.
There are eight three-dimensional irreducibles, none nontrivial on both
the $2I$ and $\SL(2,7)$ factors. This limited dimension statement does not
exclude interactions between fields in different representations.

Exact class-algebra diagonalization gives
$\Q(\sqrt2,\sqrt{-7})$ as the character field of $\SL(2,7)$, agreeing with
the independent table \cite{GroupNamesSL27}. The product character field is
\[\Q(\sqrt5,\sqrt2,\sqrt{-7}),\qquad [\Q(\sqrt5,\sqrt2,\sqrt{-7}):\Q]=8.\]
Neither $\sin(\pi/14)$ nor twice that number lies in this field.
The former seven-element split-torus assignment to principal-series
characters was incorrect.

This exclusion for one group must not be generalized to all direct
products. In $2I\times D_{14}$, where $D_{14}$ has order 14, the outer
tensor product of the defining doublet with the dihedral doublet is
irreducible and has a character value
$\varphi(\zeta_7+\zeta_7^{-1})$ of degree six. An exact resultant and
irreducibility check certify this. Factorization of an ordinary character
as a product does not make its representation reducible or prevent a
compositum of character fields.

For $N\geq3$, let $\mu=|\SL(2,\Z/N)|$. The projective index is $\mu/2$,
so $X(N)$ has $c=\mu/(2N)$ cusps and genus
$g=1+\mu/24-c/2$ \cite{SageGamma}. At level 70 this gives
$c=1728$, $g=9217$ and $\dim M_2=10944$.
The weight-two newspace dimensions for $\Gamma_0(14),\Gamma_0(35)$ and
$\Gamma_0(70)$ are respectively 1, 3 and 1. Dimension is not the number
of Galois orbits: level 35 has two orbits, one with coefficient field
$\Q(\sqrt{17})$ \cite{LMFDB35,LMFDB35Quadratic}. Modular-form coefficient fields must not
be confused with finite-group character fields. At level 70, the label
\texttt{70.2.a.a} identifies the unique rational newform orbit, not an
individual elliptic curve. LMFDB associates it with the conductor-70
elliptic-curve isogeny class \texttt{70.a} (Cremona \texttt{70a})
\cite{LMFDB70,LMFDB70Curve}. This database identification supplies no
numerical flavor fit or physical exclusion.

A zero Schur multiplier alone does not exclude central extensions with
arbitrary finite coefficients when the abelianization is nonzero. The
order-12 group
\[
\{(g,t)\in S_3\times\Z/4:\epsilon(g)=t\bmod2\}
\]
with $\epsilon(g)\in\Z/2$ the permutation parity has central kernel $\Z/2$
and does not split over $S_3$: each lift of a
transposition has order four. Likewise,
$\Delta(7350)=\Delta(150)\times_{S_3}\Delta(294)$, not their direct product.

For a twisted center $\mathcal Z(\mathrm{Vec}_G^\omega)$, the untwisted
exponent is not in general the twist-order bound. Already $G=\Z/2$ with
$\omega(a,b,c)=(-1)^{abc}$ gives projective values $\pm i$ and twist order
four. The group-theoretical Frobenius--Schur exponent instead divides
$(\exp G)^2$ \cite[Theorem 9.2]{NgSchauenburg2007}; for $2I$ the resulting
bound is 3600. This concerns canonical categorical data, not arbitrary
morphisms: in a complex-linear category the trace of
$\lambda\,\mathrm{id}_{\mathbf1}$ is $\lambda$.
Finally, the second homotopy group of a 2-group is abelian; a proposed
nonabelian $\pi_2=\SL(2,7)$ is not such a homotopy datum. A crossed module
must instead specify its boundary map, abelian kernel, action and
Postnikov class.

\section{Topological arithmetic and limits}
For a verified free finite action,
$\chi(X/G)=\chi(X)/|G|$. Integrality is necessary, not sufficient to
construct the action. Singular quotients require fixed-point data rather
than this free-cover formula. The boundary of the 600-cell triangulates
$S^3$ and has Euler characteristic zero; the filled four-ball has Euler
characteristic one. A covering deck action is an action on the covering
space, not a newly established free action on the quotient.

The relation $|\chi|/2$ counts the magnitude of a net generation index only
under the relevant Calabi--Yau threefold standard-embedding assumptions.
For example, a \emph{specified free} $\Z/2$ quotient of a $\chi=-12$ parent
would have $\chi=-6$ and standard-embedding net count three. This arithmetic
neither constructs nor excludes an asymmetric CHL model. No exhaustive
Calabi--Yau action classification or numerical flavor-fit scan is supplied.
Unsupported catalog exclusions and phenomenological fit estimates are
therefore not promoted to propositions.

\section{Conclusion}
The exact finite calculations retain substantial representation-theoretic
content while narrowing the conclusions. The literal-character arithmetic
obstruction, the $2I$ Frobenius--Schur table, the specified spinor branchings,
and the selected centralizer and congruence-group calculations are
independently testable. They do not by themselves prove a universal
chirality obstruction or rule out all finite-symmetry compactifications.
A stronger claim must state its geometric, categorical and physical inputs
and provide the corresponding proof or reproducible computation.
\bibliographystyle{plain}
\bibliography{audit_references}
\end{document}
